
设 $\mathfrak s_x$ 为扰动 $s_x$ 的多值截面。
下面说明如何在 $x^+$ 的如下仓西邻域上
诱导多值截面：
\[
(V_{x^+}^+,E_{x^+}^+,\Gamma_{x^+}^+,
\psi_{x^+}^+,s_{x^+}^+)
\]
将 $\mathfrak s_x$ 拉回到 $V_{x^+}^+$，
便在 $(W_x\times\{0\})\times V_x$ 上
得到仓西映射 $s_{x^+}^+$ 的
多值扰动的 $\pi^*E_x$ 分量
（参见 (7.3.9.5)）。
在 $(W_x\times\{0\})\times V_x$ 上，
把 $\mathbb R^{c(x^+)}$ 分量置为零。
然后延拓这个多值截面，使 (7.3.9.6) 成立。
由此得到该仓西邻域上的多值截面，
称为 $\mathfrak s_x$ 的\emph{拉回}。
\liftIndexRangeEnd{I386}

现在可以陈述证明定理~\ref{thm:7-3-1}
所需的主要技术结果。
沿用第~\ref{subsec:7-2-4} 小节的记号。
若 $P_i\in\mathcal X_{g_i}(L)$，
置 $\mathfrak d(i)=g_i$。
若 $\vec P^+=(P_1^+,\ldots,P_k^+)$，
其中 $P_i^+\in\mathcal X_{g_i^+}(L)$，
则置 $\mathfrak d^+(i)=g_i^+$。
（若 $P_j^+=L$，则 $\mathfrak d^+(j)=0$。）
照常记
$\mathcal X_{(g)}(L)=\bigcup_{g'\le g}\mathcal X_{g'}(L)$。

注意，$L$ 本身不是奇异单形。
严格说来，需要使用 $L$ 的单纯剖分。
例如，当 $k=6$、$a_1=1$、$a_2=3$、
$a_3=5$ 时，
$(L,P_1,L,P_2,L,P_3)$ 应写成
\[
\vec P^+
=\sum_{c,c',c''}\pm
(\Delta_c^n,P_1,\Delta_{c'}^n,P_2,
\Delta_{c''}^n,P_3),
\]
其中 $\Delta_c^n$ 如 (7.3.6)。
不过为简化记号，下文仍略作记号滥用，
简写为 $(L,P_1,L,P_2,L,P_3)$。

\begin{proposition}{7.3.10}\label{prop:7-3-10}
对任意 $\delta>0$、$K>0$，
存在 $L$ 上的可数链集
$\mathcal X_g(L)$（$g=1,\ldots,K$）、
当 $P_i\in\mathcal X_{(g)}(L)$、
$\|(\mathfrak d,\beta)\|\le K$ 时
$\mathcal M^{\mathrm{main}}_{k+1}(\beta;\vec P)$
上的多值截面
$\mathfrak s_{\mathfrak d,k,\beta,\vec P}$，
以及当 $P_i^+\in\mathcal X_{(g)}(L)$、
$\|(\mathfrak d^+,\beta)\|\le K$ 时
\[
[0,1]^{|\vec a|}\times
\mathcal M^{\mathrm{main}}_{k+1}(\beta;\vec P^+)
\]
上的多值截面
$\mathfrak s^+_{\vec a,\mathfrak d^+,k,\beta,\vec P^+}$，
使下列性质成立
\translatornote{源书在第二类多值截面的
输入条件中仍写 $P_i\in\mathcal X_{(g)}(L)$；
此处元组为 $\vec P^+$，依紧前定义补上 $+$。}：

\par\noindent\textup{(7.3.11.1)}\quad
$\mathcal X_{(g)}(L)$ 满足性质~\ref{prop:7-2-27}。

\par\noindent\textup{(7.3.11.2)}\quad
若 $g=\|(\mathfrak d;\beta)\|$，
则 $\mathbb Q$-链
\[
\operatorname{ev}_{0*}\left(
\mathcal M^{\mathrm{main},\mathfrak d}_{k+1}
(\beta;\vec P)
^{\mathfrak s_{\mathfrak d,k,\beta,\vec P}}
\right)
\]
可分解为 $\mathcal X_{(g)}(L)$ 中的元素。

\par\noindent\textup{(7.3.11.3)}\quad
扰动零集
$\mathcal M^{\mathrm{main},\mathfrak d}_{k+1}
(\beta;\vec P)
^{\mathfrak s_{\mathfrak d,k,\beta,\vec P}}$
位于原模空间
$\mathcal M^{\mathrm{main},\mathfrak d}_{k+1}
(\beta;\vec P)$ 的 $\delta$-邻域内。

\par\noindent\textup{(7.3.11.4)}\quad
各多值截面
$\mathfrak s_{\mathfrak d,k,\beta,\vec P}$
彼此相容，其意义由前述相容条件
\ref{cond:7-2-38}、\ref{cond:7-2-44} 给出。

\par\noindent\textup{(7.3.11.5)}\quad
若 $g=\|(\mathfrak d^+;\beta)\|$，
则 $\mathbb Q$-链
\[
\operatorname{ev}_{0*}^+\left(
\bigl([0,1]^{|\vec a|}\times
\mathcal M^{\mathrm{main},\mathfrak d^+}_{k+1}
(\beta;\vec P^+)\bigr)
^{\mathfrak s^+_{\vec a,\mathfrak d^+,k,\beta,\vec P^+}}
\right)
\]
可分解为 $\mathcal X_{(g)}(L)$ 中的元素。

\par\noindent\textup{(7.3.11.6)}\quad
对每个固定的
$(t_1,\ldots,t_{|\vec a|})\in[0,1]^{|\vec a|}$，
各多值截面
$\mathfrak s^+_{\vec a,\mathfrak d^+,k,\beta,\vec P^+}$
在
\[
\{(t_1,\ldots,t_{|\vec a|})\}\times
\mathcal M^{\mathrm{main},\mathfrak d^+}_{k+1}
(\beta;\vec P^+)
\]
上的限制，按相容条件
\ref{cond:7-2-38}、\ref{cond:7-2-44}
的意义彼此相容。

\par\noindent\textup{(7.3.11.7)}\quad
下文的相容条件 7.3.13 和 7.3.14 成立。

\par\noindent\textup{(7.3.11.8)}\quad
下文的相容条件 7.3.17 成立。

\par\noindent\textup{(7.3.11.9)}\quad
对每个 $\vec t\in[0,1]^{|\vec a|}$，
零集
\[
\bigl(\{\vec t\}\times
\mathcal M^{\mathrm{main},\mathfrak d^+}_{k+1}
(\beta;\vec P^+)\bigr)
^{\mathfrak s^+_{\vec a,\mathfrak d^+,k,\beta,\vec P^+}}
\]
在仓西邻域中与原模空间
$\mathcal M^{\mathrm{main},\mathfrak d^+}_{k+1}
(\beta;\vec P^+)$ $\delta$-接近。
\end{proposition}
